\documentclass[12pt]{article}
\usepackage[T1]{fontenc}
\usepackage[utf8]{inputenc}
\usepackage{lmodern}
\usepackage{textcomp}
\usepackage{enumitem}
\usepackage{geometry}
\usepackage{graphicx}
\usepackage{amsmath, amssymb}
\geometry{margin=1in}
\begin{document}
\begin{center}
History - Graduate Level Assessment 6 \
Total Marks: 40 \
Answer all questions.
\end{center}

\section*{Multiple Choice (10 marks)}
\begin{enumerate}[label=\arabic*.]
\item Which of the following is not a hypothesis for why hominids walked on two legs?
\begin{enumerate}[label=(\alph*)]
    \item it expended less energy than quadrupedalism for going long distances
    \item it freed the hands to carry things
    \item it allowed hominids to see greater distances
    \item it provided a more intimidating posture to potential predators
    \item it allowed for better heat regulation
\end{enumerate}
\item Which of the following crops was NOT an important component of Hopewell subsistence?
\begin{enumerate}[label=(\alph*)]
    \item squash
    \item knotweed
    \item beans
    \item sunflower
    \item maize
\end{enumerate}
\item What is one of the differences between simple foragers and complex foragers?
\begin{enumerate}[label=(\alph*)]
    \item Simple foragers employ domestication of animals.
    \item Complex foragers focus on a few highly productive resources.
    \item Complex foragers employ irrigation technology.
    \item Simple foragers use advanced hunting techniques.
    \item Complex foragers rely primarily on marine resources.
\end{enumerate}
\item How long ago did human groups begin actively controlling their food sources by artificially producing conditions under which these sources would grow?
\begin{enumerate}[label=(\alph*)]
    \item within the past 12,000 years
    \item within the past 1,000 years
    \item within the past 24,000 years
    \item within the past 2,000 years
    \item within the past 6,000 years
\end{enumerate}
\item What did the Moche build at the heart of their urban center?
\begin{enumerate}[label=(\alph*)]
    \item the Temple of the Feathered Serpent
    \item an enormous walled-in precinct of elite residences
    \item the Palace of the Painted Walls
    \item a complex irrigation system
    \item the Pyramid of the Sun
\end{enumerate}
\end{enumerate}

\section*{True / False (10 marks)}
\textit{Answer True or False}
\begin{enumerate}[label=\arabic*.]
\setcounter{enumi}{5}
\item True or False: Many non-transparent-translation theories draw on concepts from a different answer, the most obvious influence being the German theologian and philosopher Friedrich Schleiermacher.
\item True or False: In Greece, copper was known by the name chalkos (χαλκός).
\item True or False: In March, the Security Council declared a no fly zone to protect the civilian population from aerial bombardment, calling on foreign nations to enforce it; it also specifically prohibited foreign occupation. Ignoring this, Qatar sent hundreds of troops to support the dissidents, and along with France and the United Arab Emirates provided the NTC with weaponry and training.
\item True or False: The elevator motor was located at the top of the shaft or beside the bottom of the shaft.
\item True or False: However some species show circuitous migratory routes that reflect historical range expansions and are far from optimal in ecological terms. An example is the migration of continental populations of Swainson's thrush Catharus ustulatus, which fly far east across North America before turning south via Florida to reach northern South America; this route is believed to be the consequence of a range expansion that occurred about 10,000 years ago.
\end{enumerate}

\section*{Long Form Response (20 marks)}
\textit{Answer each question with a clear and complete explanation.}
\begin{enumerate}[label=\arabic*.]
\setcounter{enumi}{10}
\item Read the following passage and answer the question: Music was an important part of both secular and spiritual culture, and in the universities it made up part of the quadrivium of the liberal arts. From the early 13 th century, the dominant sacred musical form had been the motet; a composition with text in several parts. From the 1330 s and onwards, emerged the polyphonic style, which was a more complex fusion of independent voices. Polyphony had been common in the secular music of the Provençal troubadours. Many of these had fallen victim to the 13 th-century Albigensian Crusade, but their influence reached the papal court at Avignon. Question: What style of sacred musical form emerged in the 1330 s?
\item Read the following passage and answer the question: From the 1880 s to 1914, the European powers expanded their control across the African continent, competing with each other for Africa's land and resources. Great Britain controlled various colonial holdings in East Africa that spanned the length of the African continent from Egypt in the north to South Africa. The French gained major ground in West Africa, and the Portuguese held colonies in southern Africa. Germany, Italy, and Spain established a small number of colonies at various points throughout the continent, which included German East Africa (Tanganyika) and German Southwest Africa for Germany, Eritrea and Libya for Italy, and the Canary Islands and Rio de Oro in northwestern Africa for Spain. Finally, for King Leopold (ruled from 1865-1909), there was the large "piece of that grea... Question: Where in Africa did the French have control?
\end{enumerate}

\vfill
\noindent\textit{End of Paper}
\end{document}